\documentclass[a4paper,12pt]{article}

\usepackage[utf8]{inputenc}
\usepackage{geometry}
\usepackage{graphicx}
\usepackage{amsmath}
\usepackage{hyperref}

% Language and font packages
\usepackage{ctex} % Chinese support

% Pseudocode and algorithm packages
\usepackage{algorithm2e} % Advanced algorithm writing
\usepackage{algpseudocode} % Pseudocode environment for algorithms

% Fancy style packages
\usepackage{fancyhdr} % Custom headers/footers
\usepackage{color} % Color support
\usepackage{xcolor} % Extended color support
\usepackage{tcolorbox} % Colored boxes
\usepackage{framed} % Framed environments

% Math and symbol packages
\usepackage{amsfonts}
\usepackage{amssymb}
\usepackage{mathtools}

% Table and figure packages
\usepackage{booktabs} % Professional tables
\usepackage{multirow} % Multi-row tables
\usepackage{caption} % Custom captions
\usepackage{subcaption} % Subfigures

% Miscellaneous
\usepackage{enumitem} % Custom lists
\usepackage{listings} % Code listings
\usepackage{tikz} % Drawing diagrams
\usepackage{pgfplots} % For plotting graphs

\title{I have bags. Too many bags......}
\author{B13901165 Ivan Chang}
\date{\today}
\begin{document}

\maketitle

\begin{abstract}
In the textbook, we all have been forced to get familiar with the 0/1 knapsack problem.
So, why not solve a little more similar questions?
Here I prepare some related problems for you to practice.
For each problem, you are required to design the 
\begin{enumerate}
    \item DP transfer function
    \item algorithm (pseudocode)
    \item time complexity analysis
\end{enumerate}

\end{abstract}

\section{Unbounded Knapsack Problem}
\subsection*{Description}
You are given $k$ types of items. Each type $i$ has weight $w_i$ and value $v_i$,
and you may take unlimited copies of each item.
Find the maximum total value without exceeding capacity $W$.

\subsection*{Input Format}
\begin{verbatim}
k W
w1 v1
w2 v2
...
wk vk
\end{verbatim}

\subsection*{DP transfer function}
In the unbounded knapsack problem, which is actually Q9.(d) in the HW2, 
we use a 1-D array $memo[0..W]$ where $memo[j]$ represents the maximum value within capacity $j$.
for each capacity $j$, we have to iterate through all $k$ items to find the optimal solution.
The transfer function is:
\[
memo[j] = \max_{1 \leq i \leq k}(memo[j], memo[j - w_i] + v_i)
\]

\subsection*{Pseudocode}
\subsection*{Top-Down Recursive Algorithm}
\begin{algorithm}[H]
\caption{Top-Down Unbounded Knapsack}
\KwIn{$k, W, w[1..k], v[1..k]$}
\KwOut{Maximum total value}
Initialize all $memo[w] = -\infty$\;
\SetKwFunction{F}{Knapsack}
\SetKwProg{Fn}{Function}{:}{}
\Fn{\F{$w$}}{
    \If{$w = 0$}{
        \Return{$0$}\;
    }
    \If{$\text{memo}[w] \neq -\infty$}{
        \Return{$\text{memo}[w]$}\;
    }

    % try every item as a choice
    \For{$i \gets 1$ \KwTo{} $k$}{
        \If{$w_i \le w$}{
            $\text{memo}[w] \gets \max\big(\text{memo}[w],\ \text{\F}(w - w_i) + v_i\big)$\;
        }
    }

    \Return{$\text{memo}[w]$}\;
}

\Return{\text{\F}(k, W)}\;
\end{algorithm}

\subsection*{Time Complexity}
for time complexity, we have to iterate through all $W$ capacities,
and for each capacity, we iterate through all $k$ items.
Thus, the time complexity is:\[O(kW)\]
for space complexity, since we use a 1-D array $memo[0..W]$,
the space complexity is:\[O(W)\]

% ============================================
\section{Bounded Knapsack Problem}
\subsection*{Description}
You are given $k$ types of items.  
Each type $i$ has weight $w_i$, value $v_i$, and at most $n_i$ copies.  
Find the maximum total value achievable within capacity $W$.
\subsection*{Input Format}
\begin{verbatim}
k W
w1 v1 n1
w2 v2 n2
...
wk vk nk
\end{verbatim}
\subsection*{Idea}
This is the \textbf{bounded knapsack} problem.  
A naive approach is to expand all the item, and it would reduce to the 0/1 knapsack problem.
However, the time complexity would be $O(W \sum n_i)$.
To avoid expanding all $n_i$ copies, we can use \textbf{binary decomposition}:
split $n_i$ into the powers of 2.

\subsection{DP transfer function}
the transfer function is similar to the 0/1 knapsack problem:
however, we first need to do the binary decomposition to create new items.
Then, we can apply the 0/1 knapsack DP transfer function:
\[
dp[i][w] = \max(dp[i-1][w], dp[i-1][w - weight[i]] + value[i])
\]
\subsection*{Pseudocode}
\begin{algorithm}[H]
\caption{Bounded Knapsack using Binary Decomposition + 0/1 DP}
\KwIn{Capacity $W$; $k$ item types with weight $w[1..k]$, value $v[1..k]$, count $n[1..k]$}
\KwOut{Maximum total value achievable}

\BlankLine
Initialize an empty list \texttt{newItems}\;
\For{$i \gets 1$ \KwTo $k$}{
    $count \gets n[i]$, \quad $base \gets 1$\;
    \While{$count > 0$}{
        $take \gets \min(base, count)$\;
        Append item $(take \times w[i],\ take \times v[i])$ to \texttt{newItems}\;
        $count \gets count - take$\;
        $base \gets base \times 2$\;
    }
}
\BlankLine
//do the 0/1 knapsack on newItems\;
\Return $dp[W]$\;
\end{algorithm}


\subsection*{Time Complexity}
This is actually like the 0/1 knapsack problem, but with $\sum_{i=1}^k \log n_i$ items after binary decomposition.
Thus, the time complexity is:
\[O\left(W \sum_{i=1}^k \log n_i\right)\]

% ============================================
\section{Two-Knapsack Problem}
\subsection*{Description}
There are $M=2$ knapsacks, with capacities $W_1$ and $W_2$.  
You have $k$ items, each with weight $w_i$ and value $v_i$.  
Each item can be placed in \textbf{at most one} knapsack.
\subsection{Input Format}
\begin{verbatim}
k W1 W2
w1 v1
w2 v2
...
wk vk
\end{verbatim}
\subsection{DP transfer function}
for each item, we have three choices:
\begin{enumerate}
    \item not take the item
    \item put the item in the first bag
    \item put the item in the second bag
\end{enumerate}
therefore, we can build a 3-D DP array $memo[i][c1][c2]$ where \\
$i$ is the number of items, \\
$c1$ is the capacity left in the first knapsack, \\
$c2$ is the capacity left in the second knapsack.\\
the transfer function would be:
\[
memo[i][c1][c2] = \max(memo[i-1][c1][c2],\ memo[i-1][c1-w[i]][c2] + v[i],\ memo[i-1][c1][c2-w[i]] + v[i])
\]
\subsection*{Pseudocode}
\begin{algorithm}[H]
\caption{Top-Down Two-Knapsack with Memoization}
\KwIn{item count $k$; capacities $C_1, C_2$; weights $w[1..k]$; values $v[1..k]$}
\KwOut{maximum total value}
\SetKwFunction{KS}{DP}
\SetKwProg{Fn}{Function}{:}{}

\textbf{State:} $\text{memo}[0..k][0..C_1][0..C_2] \leftarrow -\infty$;

\Fn{\KS{$i, c_1, c_2$}}{
    \If{$i = 0$}{\Return $0$}
    
    \If{$\text{memo}[i][c_1][c_2] \neq -\infty$}{\Return $\text{memo}[i][c_1][c_2]$}

    % don't take item i
    $best \leftarrow$ \KS{$i-1,\, c_1,\, c_2$}\;

    % put item i into knapsack 1 (if fits)
    $best \leftarrow \max\Big(best,\ \KS(i-1,\, c_1 - w[i],\, c_2) + v[i]\Big)$\;

    $best \leftarrow \max\Big(best,\ \KS(i-1,\, c_1,\, c_2 - w[i]) + v[i]\Big)$\;

    $\text{memo}[i][c_1][c_2] = best$\;
    \Return{$\text{memo}[i][c_1][c_2]$}\;
}
\Return{\KS{$k, C_1, C_2$}}\;
\end{algorithm}


\subsection*{Time Complexity}
Since the size of the memo table is $k \times C_1 \times C_2$
the time complexity is:
\[O(k \times C_1 \times C_2)\]

% ============================================
\section{Group Knapsack Problem}
\subsection*{Description}
There are $G$ groups of items.  
Group $i$ has $t_i$ items, and each item $j$ in this group has weight $w_{ij}$ and value $v_{ij}$.  
You can choose \textbf{at most one} item from each group.
\subsection{Input Format}
\begin{verbatim}
G W
t_1
w_{11} v_{11}
w_{12} v_{12}
...
w_{1t_1} v_{1t_1}
t_2
w_{21} v_{21}
...
w_{2t_2} v_{2t_2}
...
t_G
w_{G1} v_{G1}
...
w_{Gt_G} v_{Gt_G}
\end{verbatim}
\subsection{DP transfer function}
for each group $i$, we have $t_i + 1$ choices:
\begin{enumerate}
    \item not take any item from group $i$
    \item take item $j$ from group $i$ (for each $j \in [1, t_i]$)
\end{enumerate}
therefore, we can use a 2-D DP array $memo[1...G][0..W]$ to store the maximum achievable value//
the transfer function would be:
\[
dp[i][w] = \max_{1 \leq j \leq t_i}(dp[i-1][w], dp[i-1][w - w_{ij}] + v_{ij})
\]
\subsection*{Pseudocode}
\begin{algorithm}[H]
\caption{Top-Down Group Knapsack with Memoization}
\KwIn{Number of groups $G$; capacity $W$; each group $i$ has $t_i$ items $(w_{ij}, v_{ij})$}
\KwOut{Maximum total value achievable}
\SetKwFunction{GK}{Solve}
\SetKwProg{Fn}{Function}{:}{}

\textbf{State:} $\text{memo}[0..G][0..W] = -\infty$\;

\Fn{\GK{$i, w$}}{
    \If{$i = 0$}{\Return{0}\;}
    \If{$\text{memo}[i][w] \neq -\infty$}{\Return{\text{memo}[i][w]}\;}

    $best \gets$ \GK{$i-1, w$}\;

    \For{$j \gets 1$ \KwTo $t_i$}{
        \If{$w_{ij} \leq w$}{
            $best \gets \max\big(best,\, \GK(i-1, w - w_{ij}) + v_{ij}\big)$\;
        }
    }

    $\text{memo}[i][w] \gets best$\;
    \Return{\text{memo}[i][w]}\;
}

\Return{\GK{$G, W$}}\;
\end{algorithm}

\subsection*{Time Complexity}
the size of the memo table is $G \times W$,
and for each state, we iterate through all $t_i$ items in group $i$
therefore, the time complexity is:
\[
O\left(W \sum_{i=1}^{G} t_i\right)
\]

% ============================================
\bibliographystyle{plain}
\bibliography{references}

\end{document}